\section{FREE MODULES OVER A PRINCIPAL IDEAL DOMAIN}

THEOREM 1. --- Let A be a ring such that every left ideal of A is projective (II, p.~231, Def. 1) as an A-module. Then every submodule M of a free left A-module L is a direct sum of modules isomorphic to ideals of A.

Let \((e_{\iota})_{\iota \in I}\) be a basis for L, and let \(p_{\iota}\) be the coordinate functions corresponding to this basis. Choose a well-ordering (Set Theory, III, p.~153) of I and let \(L_{\iota}\) denote the submodule generated by the \(e_{\lambda}\) for \(\lambda \leqslant \iota\) ; put \(M_{\iota} = M \cap L_{\iota}\) . The coordinate function \(p_{\iota}\) maps \(M_{\iota}\) onto an ideal \(\mathfrak{a}_{\iota}\) of A; since \(\mathfrak{a}_{\iota}\) is a projective A-module, there exists (II, p.~231, Prop. 4) a submodule \(N_{\iota}\) of \(M_{\iota}\) such that the map \(x \mapsto p_{\iota}(x)\) from \(N_{\iota}\) into \(\mathfrak{a}_{\iota}\) is bijective. Let \(M_{\iota}'\) be the submodule of L generated by the \(N_{\lambda}\) for \(\lambda \leqslant \iota\) ; we will show that \(M_{\iota}' = M_{\iota}\) for all \(\iota\) , which will imply that M is generated by the family \((N_{\iota})_{\iota \in I}\) . In fact, suppose \(M_{\lambda}' = M_{\lambda}\) for all \(\lambda < \iota\) ; then \(p_{\iota}(x) \in \mathfrak{a}_{\iota}\) for all \(x \in M_{\iota}\) ; thus there exists \(y \in N_{\iota}\) such that \(x - y\) is a linear combination of finitely many elements \(e_{\lambda}\) with \(\lambda < \iota\) ; in other words \(x - y\) is an element of \(M_{\lambda}\) for some \(\lambda < \iota\) ; the inductive hypothesis shows that \(x - y \in M_{\lambda}' \subset M_{\iota}'\) , that is that \(x \in M_{\iota}'\) , and so \(M_{\iota}' = M_{\iota}\) . It remains to show that the sum of the \(N_{\iota}\) is direct; now suppose there exists a linear relation \(\sum_{\iota} a_{\iota} = 0\) , with

\(a_{\iota} \in \mathbf{N}_{\iota}\) , where the \(a_{\iota}\) (all but finitely many of which are zero) are not all zero. Let \(\mu\) be the greatest index \(\iota\) such that \(a_{\iota} \neq 0\) ; since \(p_{\mu}(a_{\lambda}) = 0\) for \(\lambda < \mu\) , we have \(p_{\mu}(a_{\mu}) = p_{\mu}\left(\sum a_{\iota}\right) = 0\) , so \(a_{\mu} = 0\) , contradicting the choice of \(\mu\) .

COROLLARY 1. --- If every left ideal of A is projective, then every submodule of a projective left A-module is projective.

Indeed every projective A-module is a submodule of a free A-module (II, p.~231, Prop. 4), and Th. 1 applies.

COROLLARY 2. --- Every submodule of a free module over a principal ideal domain is free.

This follows immediately from Th. 1, since every ideal of a principal ideal domain is free.

COROLLARY 3. --- Every projective module over a principal ideal domain is free.

Remark. --- The proof of Th. 1 shows that every submodule of \(\mathbf{A}^{(1)}\) is isomorphic to a direct sum \(\oplus \mathfrak{a}_{\iota}\) , where each \(\mathfrak{a}_{\iota}\) is an ideal of \(\mathbf{A}\) .

PROPOSITION 1. --- If L is a free module of finite rank n over a principal ideal domain A, then every submodule M of L is a free module of rank ≤ n.

Indeed M is a free module by Cor. 2 to Th. 1, and it has rank \(\leqslant n\) by the previous remark, or by the following lemma:

Lemma 1. --- Let L be a module over a commutative ring A, generated by n elements, and let M be a free submodule of L; then M has rank ≤ n.

Suppose first that L is free. Let i denote the canonical injection of M into L. By III, p.~520, Cor., the homomorphism \(\bigwedge_{n+1}^{n+1} i: \bigwedge_{n+1}^{n+1} M \to \bigwedge_{n+1}^{n+1} L\) is injective; by III, p.~511, Prop. 6, \(\bigwedge_{n+1}^{n+1} L = \{0\}\) , so \(\bigwedge_{n+1}^{n+1} M = \{0\}\) ; it follows that M has rank \(\leqslant n\) (III, p.~518, Cor. 1). Now consider the general case; the module L is a quotient of a free module \(L'\) of rank n.~There exists a submodule \(M'\) of L isomorphic to M (II, p.~218, Prop. 21). By the first part of the argument \(M'\) has rank \(\leqslant n\) , and the result follows.

COROLLARY. --- Let E be a module over a principal ideal domain A, generated by n elements. Then any submodule F of E can be generated by at most n elements.

Indeed there exists a homomorphism \(f\) from \(\mathbf{A}^n\) onto \(\mathbf{E}\) (II, p.~216, Cor. 3), and \(f^{-1}(\mathbf{F})\) , which is a free module of rank \(m \leqslant n\) , is generated by \(n\) elements; the images of these elements under \(f\) generate \(\mathbf{F}\) .

